\Needspace{8\baselineskip}\originalchapter{2006 作业五: 磁化与磁通诊断}
题源: 17 Oct 2006, due 24 Oct. 扫描页中的积分器稳态值为\ensuremath{\pi}伏, 已人工读原图, 不用OCR的感叹号.
\begin{exercise}\label{ex:06-5-1}
原题1. 无限z pinch $\mathbf J=J(r)\mathbf e_z$, $\mathbf B=B(r)\mathbf e_\theta$, p(r), 无流无电阻, ne=ni,pe=pi,Zi1. (a)平衡p/B/J关系; (b)电子离子导引中心漂移电流JD; (c)磁化 $\mathbf M=-p\mathbf b/B$, $\mathbf J_M=\nabla\times\mathbf M$, 求JM+JD与(a)比较.
\end{exercise}
\begin{solution}
(a) $p'=-JB$, $\mu_0J=B'+B/r$. (b)中心gradB和曲率组合给 $J_D=(p/B)(1/r-B'/B)\mathbf e_z$. (c) $J_M=(1/r)(rM_\theta)'=-p'/B+pB'/B^2-p/(rB)$. 相加恰为 $-p'/B=J$. 逆磁流不能只以粒子漂移求和而遗漏闭合回旋电流, 也不能把相同压力对每个物种重复相加两次.
\end{solution}
\begin{exercise}\label{ex:06-5-2}
原题2. 10匝线圈环绕轴对称螺旋pinch, 理想积分器 $v_o=\tau^{-1}\int v_{coil}\dd t$, 已稳于\ensuremath{\pi}V. 等离子体总电流1MA, 外轴场1T. (a)τ=.1s, 等离子体对轴B小扰动, 对逆磁和顺磁分别求每米存储能; (b)n=$10^{20}$m$^{-3}$均匀、截面.25m\ensuremath{{}^{2}}, 求两情形面积平均温度.
\end{exercise}
\begin{solution}
Faraday给磁通变化模 $|\Delta\Phi|=\tau|v_o|/10=.0314159$Wb. 轴对称径向平衡乘r\ensuremath{{}^{2}}积分, 保留\ensuremath{\Delta}B一阶得 $\int p\dd A=\mu_0I^2/(8\pi)-B_0\Delta\Phi/\mu_0$. 第一项50000J/m, 第二项模25000J/m. 逆磁\ensuremath{\Delta}\ensuremath{\Phi}负, 压强积分75000J/m; 顺磁正, 25000J/m. 单原子电子离子\textbf{热内能}每米 $W_{th}=3\int p/2$ 为112500与37500J/m. 原题``plasma energy''在此诊断惯例指热能, 不把外部真空磁场全部计入; 若另求场能须给外部边界, 无限柱的外磁能发散.

两物种同温p=2nT, $\langle T\rangle=\int p\dd A/(2nA)$, 分别9362eV和3121eV. 温度以能量计, 不能把压强积分本身当单粒子T.
\end{solution}
